
\setcounter{footnote}{0}
\SGARefTarget{sga3:xxiv:chapter:xxiv}\chapter*{Exposé XXIV\\
Automorphisms of Reductive Groups}
\addcontentsline{toc}{chapter}{Exposé XXIV. Automorphisms of Reductive
Groups}
\markboth{EXPOSÉ XXIV. AUTOMORPHISMS OF REDUCTIVE GROUPS}
  {EXPOSÉ XXIV. AUTOMORPHISMS OF REDUCTIVE GROUPS}

\begin{center}
  \textit{by Michel Demazure}
\end{center}

\begingroup
\renewcommand{\thefootnote}{0}
\footnote{Editor's note: Version dated October 13, 2024.}
\addtocounter{footnote}{-1}
\endgroup
The first part of this Exposé (nos.~1 to 5) is a direct consequence of
the existence, for a reductive group, of ``sufficiently many outer
automorphisms'', a result which is itself a consequence of the weakest
form of the isomorphism theorem for pinned groups. The second part
(nos.~6 and 7) sets out two applications of the more precise results of
the preceding Exposé; in particular, no.~7 uses the generators-and-
relations theorem in its explicit form. Finally, we have included in
an appendix (no.~8) results on ``Galois'' cohomology used in the text.

Let us specify our cohomological notation. If \(S\) is a scheme and
\(G\) an \(S\)-group scheme, we write \(H^1(S,G)\) for the first
cohomology set of \(S\) with coefficients in \(G\), computed for the
\textup{(fpqc)} topology; this is also the set of isomorphism classes
of principal homogeneous \textup{(fpqc)} sheaves under \(G\). We write
\(H^1_{\mathrm{\acute et}}(S,G)\) for the corresponding set for the
étale topology; it is therefore the part of \(H^1(S,G)\) formed by
classes of homogeneous sheaves under \(G\) which are quasi-isotrivial
\textup{(that is, locally trivial for the étale topology)}. We write
\(\Fib(S,G)\) for the part of \(H^1(S,G)\) formed by classes of
representable sheaves \textup{(principal homogeneous bundles)}. Thus
we have the inclusions
\[
  H^1_{\mathrm{\acute et}}(S,G)\subset H^1(S,G),
  \qquad
  \Fib(S,G)\subset H^1(S,G).
\]

If every principal homogeneous sheaf under \(G\) is representable
\textup{(for example, if \(G\) is quasi-affine over \(S\); cf.\ SGA~1,
VIII, \SGARefLink{sga3:viii:proposition:7:9}{7.9})}, then
\[
  \Fib(S,G)=H^1(S,G).
\]

If \(S'\to S\) is a covering morphism for the \textup{(fpqc)}
topology, we write \(H^1(S'/S,G)\) for the kernel of the canonical map
\[
  H^1(S,G)\longrightarrow H^1(S',G_{S'}).
\]
It is known that \(H^1(S'/S,G)\) can be computed simplicially
\textup{(TDTE~I, \S A.4)}, which implies that, when \(S'\to S\) is a
covering for the étale topology, \(H^1(S'/S,G)\) is also the kernel of
\[
  H^1_{\mathrm{\acute et}}(S,G)
  \longrightarrow
  H^1_{\mathrm{\acute et}}(S',G_{S'}).
\]

Finally, following \SGARefLink{sga3:viii:corollary:4:5}{Exp.~VIII, 4.5}, the name ``\SGARefLink{sga3:viii:corollary:4:5}{Theorem~90}'' is given to
the following assertion: ``every principal homogeneous sheaf under
\(G_{m,S}\) is representable and locally trivial.'' This assertion is
equivalent to
\[
  H^1(S,G_{m,S})=\operatorname{Pic}(S),
\]
or again to
\[
  H^1(S,G_{m,S})=0
\]
when \(S\) is local \textup{(or, more generally, semilocal)}.

\SGARefTarget{sga3:xxiv:section:1}
\section*{1. Scheme of Automorphisms of a Reductive Group}
\addcontentsline{toc}{section}{1. Scheme of Automorphisms of a Reductive Group}

\SGARefTarget{sga3:xxiv:1.0}
\subsection*{1.0.}
We must first make precise certain definitions from the preceding
Exposé. Let \(S\) be a nonempty scheme, let \(G\) be a reductive
\(S\)-group, and let
\[
  \mathscr R=(M,M^*,R,R^*,\Delta)
\]
be a pinned reduced root datum \textup{(\SGARefLink{sga3:XXIII:1.5}{Exp.~XXIII, 1.5})}. A
\emph{pinning of \(G\) of type \(\mathscr R\)}, or an
\(\mathscr R\)-\emph{pinning of \(G\)}, consists of:
\begin{enumerate}
\item[(i)] an isomorphism of \(D_S(M)\) onto a maximal torus \(T\) of
  \(G\) \textup{(or, equivalently, a monomorphism
  \(D_S(M)\to G\) whose image is a maximal torus \(T\) of \(G\))},
  identifying \(R\) with a root system of \(G\) relative to \(T\)
  \textup{(\SGARefLink{sga3:xix:definition:3-6}{Exp.~XIX, 3.6})} and \(R^*\) with the corresponding set of
  coroots;
\item[(ii)] for each \(\alpha\in\Delta\), an element
  \[
    X_\alpha\in\Gamma(S,\mathfrak g^\alpha)^\times.
  \]
\end{enumerate}
The group \(G\) possesses an \(\mathscr R\)-pinning if and only if it
is split and of type \(\mathscr R\) \textup{(\SGARefLink{sga3:xxii:def:2.7}{Exp.~XXII, 2.7})}.

If \(u:G\to G'\) is an isomorphism of reductive \(S\)-groups, then
every \(\mathscr R\)-pinning \(\mathcal E\) of \(G\) determines, by
``transport of structure'', an \(\mathscr R\)-pinning
\(u(\mathcal E)\) of \(G'\). If
\(v:\mathscr R'\to\mathscr R\) is an isomorphism of pinned root data,
then every \(\mathscr R'\)-pinning \(\mathcal E'\) of \(G'\)
determines, by transport of structure, an \(\mathscr R\)-pinning
\(v(\mathcal E')\) of \(G\).

Call a \emph{pinned group} a triple \((G,\mathscr R,\mathcal E)\),
where \(G\) is a reductive \(S\)-group, \(\mathscr R\) is a pinned
reduced root datum, and \(\mathcal E\) is a pinning of \(G\) of type
\(\mathscr R\). An isomorphism from the pinned group
\((G,\mathscr R,\mathcal E)\) onto the pinned group
\((G',\mathscr R',\mathcal E')\) is a pair \((u,v)\), where
\(u:G\to G'\) is an isomorphism and
\(v:\mathscr R'\to\mathscr R\) is an isomorphism of pinned root data,
such that
\[
  u(\mathcal E)=v(\mathcal E').
\]
\footnote{Editor's note: Thus
\(\Lie(u)(X_{v(\alpha)})=X'_\alpha\) for every
\(\alpha\in R'\).}

\textit{N.B.} If \(S\) is nonempty, \(v\) is uniquely determined by
\(u\), and by an abuse of language \(u\) will also be called an
isomorphism of the pinned groups. In particular, if
\((G,\mathscr R,\mathcal E)\) is a pinned group, an automorphism of
\((G,\mathscr R,\mathcal E)\) is therefore an automorphism \(u\) of
\(G\) for which there exists an automorphism \(v\) of \(\mathscr R\)
such that \(u(\mathcal E)=v(\mathcal E)\). It is thus an automorphism
of \(G\) which normalizes \(T\), induces on \(T\) an automorphism of
the form \(D_S(h)\), where \(h\) is an automorphism of \(M\), and
permutes the elements \(X_\alpha\), \(\alpha\in\Delta\).
\textup{(As is readily seen, the preceding conditions in fact
characterize the automorphisms of
\((G,\mathscr R,\mathcal E)\).)}

There is an evident contravariant functor
\[
  \Phi:\quad
  (G,\mathscr R,\mathcal E)\longmapsto\mathscr R,
  \qquad
  (u,v)\longmapsto v,
\]
and the principal result of the preceding Exposé
\textup{(\SGARefLink{sga3:xxiii:theorem:4-1}{Exp.~XXIII, 4.1})} shows that this functor is fully faithful
\textup{(we shall moreover see in the following Exposé that it is an
equivalence of categories)}. In particular, it follows that the
automorphism group of \((G,\mathscr R,\mathcal E)\) is canonically
isomorphic to the automorphism group of the pinned root datum
\(\mathscr R\) \textup{(cf.\ \SGARefLink{sga3:XXIII:5.5}{Exp.~XXIII, 5.5})}.

\SGARefTarget{sga3:xxiv:1.1}
\subsection*{1.1.}
Let \(S\) be a scheme; endow \((\mathbf{Sch})/S\) with the
\textup{(fpqc)} topology, and consider the \(S\)-sheaf of groups
\(\underline{\Aut}_{S\text{-}\mathrm{gr.}}(G)\), where \(G\) is an
\(S\)-group scheme. There is an exact sequence of \(S\)-sheaves of
groups
\[
  1\longrightarrow \underline{\Centr}(G)
  \longrightarrow G
  \xrightarrow{\mathrm{int}}
  \underline{\Aut}_{S\text{-}\mathrm{gr.}}(G),
\]
which defines a monomorphism
\[
  j:G/\underline{\Centr}(G)
  \longrightarrow \underline{\Aut}_{S\text{-}\mathrm{gr.}}(G).
\]
The image sheaf of \(j\) is the sheaf of inner automorphisms of \(G\).
An automorphism \(u\) of \(G\) is inner if and only if there is a
covering family \(\{S_i\to S\}\) and, for every \(i\), an element
\(g_i\in G(S_i)\) such that
\[
  \operatorname{int}(g_i)=u_{S_i}.
\]
In this case, if \(v\) is another automorphism of \(G\), one sees at
once that
\[
  \operatorname{int}(v)u=vuv^{-1}
\]
is the inner automorphism defined by the family
\(g'_i=v(g_i)\). It follows that the image of \(j\) is normal in
\(\underline{\Aut}_{S\text{-}\mathrm{gr.}}(G)\). The quotient sheaf of
groups, denoted \(\underline{\Autext}(G)\), is the sheaf of outer
automorphisms of \(G\). Thus there is an exact sequence
\[
  1\longrightarrow G/\underline{\Centr}(G)
  \longrightarrow \underline{\Aut}_{S\text{-}\mathrm{gr.}}(G)
  \longrightarrow \underline{\Autext}(G)
  \longrightarrow 1.
\]
All the preceding definitions are compatible with base change. They
are naturally valid in any site.

\SGARefTarget{sga3:xxiv:1.2}
\subsection*{1.2.}
Let \(S\) be a scheme and let \((G,\mathscr R,\mathcal E)\) be a
pinned reductive \(S\)-group. Let \(E\) be the \textup{(abstract)}
automorphism group of the pinned root datum \(\mathscr R\), that is,
the group of automorphisms of \(\mathscr R\) which normalize
\(\Delta\). By \SGARefLink{sga3:XXIII:5.5}{Exp.~XXIII, 5.5}, there is a canonical monomorphism
\[
  E\longrightarrow \underline{\Aut}_{S\text{-}\mathrm{gr.}}(G)
\]
which associates with \(h\in E\) the unique automorphism \(u\) of the
pinned group \(G\) such that \(\Phi(u)=h\). This monomorphism
canonically defines a monomorphism of sheaves
\[
  \tag{*}\label{sga3:xxiv:subsection:1-2:equation:star}
  a:E_S\longrightarrow \underline{\Aut}_{S\text{-}\mathrm{gr.}}(G).
\]
An automorphism \(u\) of \(G\) is a section of the image sheaf of \(a\)
if and only if the following conditions hold:
\begin{enumerate}
\item[(i)] \(u\) normalizes \(T\). It is then known that \(u\)
  permutes the roots of \(G\) relative to \(T\). If \(\alpha\in R\),
  then
  \[
    u(\alpha):t\longmapsto\alpha\bigl(u^{-1}(t)\bigr)
  \]
  is therefore locally on \(S\) of the form
  \(t\longmapsto\beta(t)\), with \(\beta\in R\). The second condition
  may then be written as follows:
\item[(ii)] If \(\alpha\in\Delta\) and \(U\) is an open subset of
  \(S\) such that \(u(\alpha)_U\in R\), then
  \(u(\alpha)_U\in\Delta\) and
  \[
    \Lie(u_U)(X_\alpha)_U=(X_{u(\alpha)})_U.
  \]
\end{enumerate}
It follows at once from the definitions that the sections of
\(a(E_S)\) normalize the subgroups of \(G\) defined by the pinning:
\[
  T,\quad B,\quad B^{-},\quad U,\quad U^{-}.
\]

With these definitions in place, we have:
